\section{Теорема о замкнутом графике.}
\begin{notet}
    Эта секция и следующая не были прочитаны на лекциях, но входят в программу досрочного экзамена.
\end{notet}
\begin{definition}
    Будем называть графиком оператора $A\colon X \to Y$ и обозначать его как $\graph A$ множество
    \[
        \graph A = \{(x, Ax) \mid x \in X\} \subset X \times Y.
    \]
\end{definition}
\begin{theorem}[о замкнутом графике]
    Для оператора $A$ между банаховыми пространствами $X$ и $Y$ следующие эквивалентны:
    \begin{enumerate}
        \item $A$ -- непрерывен;
        \item $\graph A$ замкнут в $X \times Y$.
    \end{enumerate}
\end{theorem}
\begin{proof}
    Пусть $A$ -- непрерывен и $\{(x_n, y_n)\} \subset \graph A$, причём $y_n = Ax_n$.
    Если $(x_n, y_n) \rightarrow (x, y) \in X \times Y$, то, в силу непрерывности оператора,
    \[
        y = \lim y_n = \lim A x_n = A \lim x_n = Ax.
    \]
    И, следовательно, $(x, y) \in \graph A$.
    Таким образом $\graph A$ -- замкнут.

    Пусть $\graph A$ -- замкнут в $X \times Y$.
    Заметим, что $\graph A$ -- замкнутое линейное подпространство $X \times Y$, а следовательно оно является и банаховым.
    Рассмотрим каноническую проекцию $\pi_X\colon (x, Ax) \mapsto x$.
    Заметим, что $\pi_X$ -- непрерывная биекция между банаховыми пространствами ($\graph A$ и $X$, соответственно), а следовательно непрерывно обратимо по теореме об обратном операторе.
    Таким образом
    \[
        A = \pi_Y \circ \pi^{-1}_X: x \xrightarrow{\pi^{-1}_X} (x, Ax) \xrightarrow{\pi_Y} Ax.
    \]
    и оператор $A$ -- непрерывен как композиция непрерывных отображений.
\end{proof}
\section{Теорема Хеллингера-Теплица.}
\begin{theorem}[Хеллингер, Теплиц]
    Пусть $H$ -- гильбертово пространство, $A$ -- симметричный линейный оператор на нём, то есть
    \[
        \forall x, y \in H \hookrightarrow \langle x, Ay \rangle = \langle Ax, y \rangle.
    \]
    Тогда $A$ -- непрерывен.
\end{theorem}
\begin{proof}
    Пусть $x_k \rightarrow x_0$ и $Ax_k \rightarrow y_0$.
    Тогда
    \[
        \langle Ax_{k}, z \rangle \rightarrow \langle y_0, z \rangle.
    \]
    С другой стороны
    \[
        \langle Ax_{k}, z \rangle = \langle x_{k}, Az \rangle \rightarrow \langle x_0, Az \rangle = \langle Ax_0, z \rangle.
    \]
    Поскольку это верно для произвольного $z \in H$, то $Ax_0 = y_0$ и график $A$ -- замкнут.
    По теореме о замкнутом графике $A$ -- непрерывен.
\end{proof}